% Vetores unitários em coordenadas cartesianas
% Feito para o ÉguaTeX (eguatex.com.br). Uso livre (CC0 1.0).
%
% Compile e pronto: sai só a figura, em PDF.
% Para pôr no seu trabalho, use o PDF com \includegraphics ou copie
% o bloco tikzpicture para dentro de uma figure.
\documentclass[tikz,border=4pt]{standalone}
\usepackage[T1]{fontenc}
\usepackage{lmodern}
\usepackage{amsmath}

\colorlet{corVet}{red!75!black}

\begin{document}
\begin{tikzpicture}[
  x={(-0.4435cm,-0.3663cm)}, y={(0.9768cm,-0.1715cm)}, z={(0cm,1cm)},
  line cap=round, line join=round, >=stealth,
  eixo/.style={line width=0.75pt},
  guia/.style={line width=0.55pt, black!35, dashed},
  vetor/.style={->, line width=1.4pt, corVet}]

% Mude estes números e o desenho e os rótulos acompanham.
\def\px{2.5} \def\py{3.2} \def\pz{2.4} % ponto P = (x, y, z)
\def\tam{1.3}                          % tamanho dos vetores, em cm

% eixos
\draw[eixo] (0,0,0) -- (5,0,0) node[below left] {$x$};
\draw[eixo] (0,0,0) -- (0,5.5,0) node[right] {$y$};
\draw[eixo] (0,0,0) -- (0,0,4.8) node[above] {$z$};

% caixa de guias até o ponto
\draw[guia] (\px,0,0) -- (\px,\py,0) -- (0,\py,0);
\draw[guia] (\px,\py,0) -- (\px,\py,\pz);
\fill (\px,\py,\pz) circle[radius=1.8pt] node[left=2pt] {$P$};

% vetores unitários no ponto P
\draw[vetor] (\px,\py,\pz) -- ++(\tam,0,0) node[below left=-1pt] {$\hat{\mathbf{x}}$};
\draw[vetor] (\px,\py,\pz) -- ++(0,\tam,0) node[right] {$\hat{\mathbf{y}}$};
\draw[vetor] (\px,\py,\pz) -- ++(0,0,\tam) node[above] {$\hat{\mathbf{z}}$};
\end{tikzpicture}
\end{document}
